% !TEX root = ../../../proposal.tex

\section{Diffusion Models}
\label{sec:diffusion_models}

\subsection{Denoising Diffusion Probabilistic Models}

Diffusion models~\citep{pmlr-v37-sohl-dickstein15, Ho2020} define a two-process framework
for generative modelling. The \emph{forward process} gradually corrupts a data sample
$x_0 \sim q(x_0)$ by adding Gaussian noise over $T$ discrete steps:
\begin{equation}
    q(x_t \mid x_{t-1}) = \mathcal{N}\!\left(x_t;\,\sqrt{1-\beta_t}\,x_{t-1},\,\beta_t I\right),
\end{equation}
where $\beta_1, \ldots, \beta_T$ is a fixed variance schedule. Marginalising over
intermediate steps gives a closed-form expression for $x_t$ given $x_0$ directly:
\begin{equation}
    q(x_t \mid x_0) = \mathcal{N}\!\left(x_t;\,\sqrt{\bar\alpha_t}\,x_0,\,(1-\bar\alpha_t)I\right),
    \qquad \bar\alpha_t = \prod_{s=1}^t (1-\beta_s),
\end{equation}
so that $x_t = \sqrt{\bar\alpha_t}\,x_0 + \sqrt{1-\bar\alpha_t}\,\varepsilon$ for
$\varepsilon \sim \mathcal{N}(0,I)$. As $t \to T$, $\bar\alpha_t \to 0$ and $x_T$ is
approximately isotropic Gaussian noise. The \emph{reverse process} learns to invert this
corruption:
\begin{equation}
    p_\theta(x_{t-1} \mid x_t) = \mathcal{N}\!\left(x_{t-1};\,\mu_\theta(x_t, t),\,\sigma_t^2 I\right).
\end{equation}
Rather than predicting $\mu_\theta$ directly, Ho et al.~\citep{Ho2020} parameterise the
model as a noise predictor $\varepsilon_\theta(x_t, t)$ and train with the simplified
objective
\begin{equation}
    \mathcal{L}_{\mathrm{simple}} = \mathbb{E}_{t,\,x_0,\,\varepsilon}\!\left[
    \|\varepsilon - \varepsilon_\theta(x_t, t)\|^2\right].
\end{equation}
Sampling proceeds by drawing $x_T \sim \mathcal{N}(0,I)$ and iteratively denoising through
$x_T \to x_{T-1} \to \cdots \to x_0$ using $\varepsilon_\theta$.

\subsection{Score-Based Formulation}

Song et al.~\citep{song2021scorebased} reframed diffusion models in continuous time via
stochastic differential equations. The forward SDE takes the form $\mathrm{d}x =
f(x,t)\,\mathrm{d}t + g(t)\,\mathrm{d}w$, and the corresponding reverse SDE requires
the \emph{score function} $\nabla_x \log p_t(x)$. A score network $s_\theta(x,t) \approx
\nabla_x \log p_t(x)$ is learned by denoising score matching, and it is connected to the
DDPM noise predictor by
\begin{equation}
    s_\theta(x_t, t) = -\frac{\varepsilon_\theta(x_t, t)}{\sqrt{1-\bar\alpha_t}}.
    \label{eq:score_noise}
\end{equation}
This score-function perspective is central to Phase~1: all compositional operators studied
there act by combining score functions from independently conditioned model passes, rather
than by modifying model weights or the decoding architecture.

\subsection{Conditional Generation and Classifier-Free Guidance}

To condition generation on a context $c$ (a text prompt, class label, or embedding), the
model learns a conditional noise predictor $\varepsilon_\theta(x_t, t, c)$ jointly with an
unconditional predictor $\varepsilon_\theta(x_t, t, \varnothing)$ via random conditioning
dropout. Classifier-free guidance (CFG)~\citep{ho2022classifier} amplifies the conditional
direction in score space at inference time:
\begin{equation}
    \tilde\varepsilon(x_t, t, c)
    = \varepsilon_\theta(x_t, t, \varnothing)
    + w\bigl[\varepsilon_\theta(x_t, t, c) - \varepsilon_\theta(x_t, t, \varnothing)\bigr],
    \label{eq:cfg}
\end{equation}
where $w > 1$ is the guidance weight. Increasing $w$ sharpens the sample toward $c$ at the
cost of diversity. In Phase~1, single-prompt CFG serves as the \emph{semantic composition}
baseline --- the monolithic joint prompt $c = \text{``a cat and a dog''}$ --- against which
product-of-experts logical composition is compared. In Phase~3, per-disease CFG conditioning
provides the individual score fields that are combined at inference time to synthesise
co-morbid chest X-rays.
